% profiles/university-of-south-florida.tex — University of South Florida
% ============================================================================
% source: https://www.usf.edu/graduate-studies/students/electronic-thesis-dissertation/etd-formatting-requirements/general/index.aspx
%         (Office of Graduate Studies, ETD Formatting Requirements → General →
%         Page Margins)
% source: .../etd-formatting-requirements/general/line-spacing.aspx
% source: .../etd-formatting-requirements/general/font-options.aspx
% source: .../etd-formatting-requirements/general/organization.aspx
% source: .../etd-formatting-requirements/section-specific/index.aspx
%         (Overview, and "Key Requirements for ADA Compliant PDFs")
% source: .../etd-formatting-requirements/section-specific/title-page.aspx
%         (the title-page specification, which is prescriptive down to the
%         number of blank lines between blocks)
% accessed: 2026-09-12
% checked:  texlib-thesis-profile-round
%
% VERIFIED and set: geometry, spacing, the title page, the absence of an
% approval page, tagging, and the accessibility requirements. Nothing is left to
% the neutral fallback.
%
% USF IS UNUSUALLY BLUNT ABOUT LaTeX, and anyone filing here with this class
% should read it first, quoted in full:
%
%   "Due to the difficulty creating Accessible PDFs through LaTex, we are
%    currently discouraging the use of LaTex. Please note that many of the
%    resources and supports availble for LaTex are limited or out of date. USF
%    does not have any official LaTex support or an official LaTex Template.
%    Please use an outside source to check your PDF for accessibilty before
%    submitting to USF. Please do not use Adobe's auto tagging feature as it mis
%    identifies many aspects of the navigational structureand is particularly
%    problematic for equations inappropriatly tagged as text. For this reason, we
%    encourage the use of Word instead." [sic throughout]
%
% and, in the ADA section: "In LaTex there are packages that allow for the
% document structure to be read and converted into a tagged PDF when compiled."
% That is what this class does, and USF names veraPDF as an acceptable checker —
% "We recommend using PAC or VeraPDF to confirm your PDF is fully accessible
% before submitting to ProQuest" — which is the same tool the TeXLib build gate
% runs. A tagged build from this class is therefore aimed squarely at what USF
% asks for, but the discouragement is USF's own position and is recorded here
% rather than argued with.

\thesisinstitution{University of South Florida}

% "Page margins must be set to 1\" on each side throughout the entire
% manuscript. All sections of the manuscript have the same margin requirements."
% and "Ensure that content (tables, figures, text) does not overflow into the
% margins." No binding allowance is offered and none is implied — submission is
% electronic, through ProQuest.
\thesissetgeometry{left=1in, right=1in, top=1in, bottom=1in}

% "The majority of your thesis/dissertation should be double-spaced", with the
% dedication, acknowledgments, abstract and body listed as double and the front
% matter lists as single. The title page is single-spaced and sets its own
% spacing below. \thesissetchapteropening is NOT called: USF states no
% chapter-opening exception, and in fact "does not require a specific structure
% or organization of the content in the main body".
\thesissetspacing{double}

% NOT set, for want of a hook, but worth knowing: "Choose one primary font size,
% 10 pt., 11 pt., or 12 pt., and use this one chosen size throughout for main
% body text, headings, and page numbers." The typeface rule ends with an
% explicit permission this class happens to satisfy by default: "Students using
% LaTeX may use Computer Modern Roman."

% --- Accessibility -----------------------------------------------------------
% USF's ADA section is a requirement, not advice: "Please check your PDF for
% accessibility compliance prior to submission", and the Overview warns that
% students must "ensure that the headings are bookmarked, accurate, and that the
% PDF is ADA compliant".
%
% \thesisrequirepdfstandard is NOT called. PDF/UA appears only inside the name of
% a tool USF recommends ("We recommend using PAC or VeraPDF"), never as a
% standard the document must declare. \thesisrequiredocumentlanguage is NOT
% called either: no USF page asks for the document language, even though it
% names WCAG 2.1 AA for contrast.
%
% USF's WCAG citation is scoped to colour contrast specifically, and it carries
% no ADA compliance date at all. Nothing here to reconcile against another
% profile's date.
\thesisrequiretagging
\thesisaccessibilityrequirement{Tagging and structure}
	{"Use Styles in Word to define headings (Heading 1, Heading 3,...) and tag
	 the Title (as Title) ... so that screen readers can interpret the document
	 structure." The title must be tagged as Title, "different from Heading 1".
	 "Do not use Adobe's Auto Tagging feature, it erroneously tags structural
	 elements and is not ADA compliant as of yet."}
\thesisaccessibilityrequirement{Bookmarks}
	{"all level headings must be bookmarked in the PDF to reflect the
	 organization of the manuscript" — "PDFs submitted to ETD must be
	 bookmarked, as this brings the organizational structure from the Word
	 document into the PDF".}
\thesisaccessibilityrequirement{Searchable text}
	{"PDFs must contain actual text, not just images of text. Tables must be
	 actual text and not images of text. Formulas and equations should be
	 inserted as such, and searchable."}
\thesisaccessibilityrequirement{Alt text}
	{"Please make sure all figures have accurate Alt Text present to ensure that
	 your PDF is ADA Compliant ... One to two sentences that describe the figure
	 is all that is needed."}
\thesisaccessibilityrequirement{Reading order}
	{"Ensure the document flows in a logical sequence for assistive technology.
	 Follow the organizational structure required by USF's Institutional
	 Guidelines."}
\thesisaccessibilityrequirement{Colour contrast}
	{"Text must meet WCAG 2.1 AA color contrast ratios. For USF, black font on a
	 white background must be used for body text, headings, and Figure/Table
	 titles and captions." And: "When using color to convey information, do not
	 rely on color alone."}
\thesisaccessibilityrequirement{Checking}
	{"Please check your PDF for accessibility compliance prior to submission. We
	 recommend using PAC or VeraPDF to confirm your PDF is fully accessible
	 before submitting to ProQuest."}

% --- Title page --------------------------------------------------------------
% "The Title Page format is very specific ... Even the amount of space between
% text is specified." USF means it: the page below is that specification, block
% by block, with its blank-line counts reproduced as multiples of \baselineskip.
%
%   Line spacing: single. Page number: none. Alignment: every line centred.
%   Title starts 2in from the top edge, "(1" below the 1" margin)", in Title
%   Case, "no special text format such as Bold or Italics", faux double-spaced
%   if it runs to a second line, and no more than two lines.
%   Then 3 blank lines / "by" / 3 blank lines / the author's full name, with
%   "no credentials" / 4 blank lines.
%   Then the degree block, broken exactly as USF breaks it, ending "University
%   of South Florida".
%   Then 3 blank lines — or 2 if the optional concentration line is present,
%   which this file switches automatically.
%   Then the committee, "One professor per line", the Major Professor first and
%   the only one carrying a title.
%   Then 2 blank lines / "Date of Approval:" / the defence date.
%   Then 3 blank lines / the keywords / 1 blank line / the copyright line.
%
% FIVE THINGS A FILER MUST SET. USF's title page names more fields than the
% class's metadata model holds, so the profile provides them; each is empty by
% default and an empty one simply drops its line.
%   \usfdepartment{...}  — "Confirm your proper Department name is listed
%                          correctly and spelled out completely." Defaults to
%                          \thesis@program, which is often right but is not the
%                          same field.
%   \usfcollege{...}     — "Make sure the College is listed correctly and
%                          spelled out."
%   \usfconcentration{...} — optional; printed as "with a concentration in X",
%                          and its presence is what changes the following gap
%                          from 3 blank lines to 2.
%   \usfkeywords{...}    — "4-6 words or phrases not used in the Title ...
%                          separated by commas."
%   \usfcopyrightyear{...} — USF uses TWO dates on this page: the Date of
%                          Approval is "the date the thesis or dissertation was
%                          successfully defended" (that is \graddate, formatted
%                          "Month Day, Year"), while the copyright year is "the
%                          year matching the year of the final submission".
%                          Defaults to \thesis@date so an unset value shows up
%                          as a full date where a bare year belongs, rather than
%                          quietly printing the wrong one.
%
% The committee: "Only the Major Professor or Co-Major Professors have titles
% before their names", and everyone else is "Name, Degree Credentials". So the
% Major Professor line is built from \gradadvisor and the remaining members from
% \committeemember — do NOT also list the major professor as a \committeemember,
% or they appear twice. Roles passed to \committeemember are dropped here,
% because USF lists none.
\newcommand{\thesis@usf@dept}{\thesis@program}
\newcommand{\usfdepartment}[1]{\renewcommand{\thesis@usf@dept}{#1}}
\newcommand{\thesis@usf@college}{}
\newcommand{\usfcollege}[1]{\renewcommand{\thesis@usf@college}{#1}}
\newcommand{\thesis@usf@conc}{}
\newcommand{\usfconcentration}[1]{\renewcommand{\thesis@usf@conc}{#1}}
\newcommand{\thesis@usf@keywords}{}
\newcommand{\usfkeywords}[1]{\renewcommand{\thesis@usf@keywords}{#1}}
\newcommand{\thesis@usf@cryear}{\thesis@date}
\newcommand{\usfcopyrightyear}[1]{\renewcommand{\thesis@usf@cryear}{#1}}

\renewcommand{\thesistitlepage}{%
	\loadgeometry{nopagenumber}%
	\begin{titlepage}%
		\thispagestyle{empty}\singlespacing\centering
		% "Title should start 2\" from the top edge of the page (1in below the
		% 1in margin)."
		\vspace*{1in}%
		% "Title is faux double-spaced if it extends to more than one line."
		\begingroup\doublespacing
		\thesis@tagtitle{Title}{\@title}\par
		\endgroup
		\vspace{3\baselineskip}
		by\par
		\vspace{3\baselineskip}
		\@author\par
		\vspace{4\baselineskip}
		A \thesis@docnoun{} submitted in partial fulfillment\\
		of the requirements for the degree of\\
		\thesis@degree\\
		\ifx\thesis@usf@conc\@empty\else with a concentration in \thesis@usf@conc\\\fi
		\thesis@usf@dept\\
		\ifx\thesis@usf@college\@empty\else \thesis@usf@college\\\fi
		University of South Florida\par
		% "Add 3 single-spaced lines ... if no concentration is listed"; "Add 2
		% ... if the (optional) concentration line is listed".
		\ifx\thesis@usf@conc\@empty
			\vspace{3\baselineskip}
		\else
			\vspace{2\baselineskip}
		\fi
		\begingroup
		% "One professor per line", no roles after the names.
		\renewcommand{\thesis@memberline}[2]{##1\par}%
		Major Professor: \thesis@advisor\par
		\thesis@committee
		\endgroup
		\vspace{2\baselineskip}
		Date of Approval:\\
		\thesis@date\par
		\vspace{3\baselineskip}
		\ifx\thesis@usf@keywords\@empty\else
			Keywords: \thesis@usf@keywords\par
			\vspace{\baselineskip}
		\fi
		Copyright \textcopyright\ \thesis@usf@cryear, \@author\par
		\vfill
	\end{titlepage}%
	\loadgeometry{normal}%
}

% --- Committee approval page --------------------------------------------------
% USF prescribes none, and the committee is named on the title page instead —
% "(Co-)Major Professor(s) and committee members' full names and degree
% credentials are listed" there, together with the "Date of Approval", which is
% "the date the thesis or dissertation was successfully defended".
%
% The Office of Graduate Studies enumerates the sections of an ETD one page at a
% time — Title Page, Table of Contents, List of Tables, List of Figures,
% Abstract, Main Body, Tables, Figures, References, Optional Sections — and no
% approval, committee or signature page appears among them, nor among the
% optional sections.
%
% This rests on an absence rather than a sentence, and is flagged as such for the
% reviewer: deleting this line restores the class's neutral page.
\thesisnoapprovalpage
